\documentclass[11pt]{article}
\usepackage[a4paper,margin=20mm]{geometry}
\usepackage{amsmath,amssymb,amsthm,microtype}
\usepackage[colorlinks=true,linkcolor=blue,urlcolor=blue]{hyperref}
\newtheorem{theorem}{Theorem}
\newtheorem{lemma}[theorem]{Lemma}
\newtheorem{corollary}[theorem]{Corollary}
\newtheorem{remark}[theorem]{Remark}
\title{Circuits for permutation MPUs and finite-phase monomial MPUs}
\author{TNLean contributors}
\date{2 October 2026}
\begin{document}
\maketitle

\section{Statement and scope}

Let $U$ act on $N\geq1$ qubits. Write its coefficients as a function of the
paired physical indices, with output listed first:
\begin{align}
 f(a_1,b_1;\ldots;a_N,b_N)
 &=\langle a_1\cdots a_N|U|b_1\cdots b_N\rangle.
 \label{eq:coefficients}
\end{align}
For $0\leq k\leq N$, let $M_k$ be the matrix with rows indexed by the
first $k$ pairs and columns by the remaining $N-k$ pairs. An open-boundary
MPO of bond dimension at most $D$ gives a factorization of $M_k$ through a
vector space of dimension at most $D$. Consequently $\operatorname{rank}M_k\leq D$.
This rank is the operator Schmidt rank at the corresponding cut, rather
than the rank of $U$ as an operator on the full chain.

\begin{theorem}\label{thm:monomial}
Suppose $U$ is monomial: there are a permutation $F$ of $\{0,1\}^N$ and
phases $z(x)$ of modulus one such that $U|x\rangle=z(x)|F(x)\rangle$.
Suppose the set of values of $z$ has at most $m$ elements, and
$\operatorname{rank}M_k\leq D$ at every cut, where $D\geq1$.
Put $W=(m+1)^D$. Then $U$ has an exact circuit with arbitrary one- and
two-qubit gates, using
\begin{align}
 \text{gate count}&=O(NW^3),&
 \text{auxiliary qubits}&=O(NW^3),&
 \text{depth}&=O\bigl(\log(N+1)\log(W+1)\bigr).
 \label{eq:resources}
\end{align}
All auxiliary qubits begin and end in $|0\rangle$.
The constants in these estimates are absolute.
\end{theorem}

\begin{corollary}\label{cor:permutation}
Every permutation MPU with maximal bond dimension $D$ has an exact circuit
of $O(N8^D)$ gates and auxiliary qubits, and depth
$O(D\log(N+1))$, with arbitrary two-qubit interactions.
In particular, fixed bond dimension gives linear circuit size and
logarithmic depth, without a lower bound on Schmidt coefficients.
More generally, $D\leq c\log_2(N+1)$ for a fixed $c$ gives gate count and
auxiliary space $O(N(N+1)^{3c})$ and depth $O(\log^2(N+1))$.
\end{corollary}

The finite phase alphabet restricts monomial MPUs, but imposes no further
hypothesis on permutation MPUs, whose nonzero entries are all one.
The dependence on $D$ is exponential. Thus this is a polynomial bound in
$N$ for fixed $D$, not a polynomial bound in $N,D$ for variable $D$.
The logarithmic depth permits arbitrary pairs of qubits to interact. It is
not a bound for geometrically local circuits. Routing every gate on a line
by swaps gives the weaker, but still polynomial for fixed $D$, bound
$O(N^2W^6)$ on gate count and depth.

Styliaris--Trivedi--Cirac's general construction depends on conditioning
\cite{STC25}; see \nolinkurl{references/2508.08160/main.tex}, lines 1373--1402.
The argument here uses finite coefficient values and leaves general MPUs
unresolved. No claim of priority for the underlying finite-state or
reversible-computation arguments is intended.

\section{Rank bounds the number of residual functions}

\begin{lemma}\label{lem:rows}
Let $M$ be a matrix over a field, with finitely many columns and rank $r$.
If its entries belong to a finite set $S$ of cardinality $s$, then $M$ has
at most $s^r$ distinct rows.
\end{lemma}

\begin{proof}
Choose $r$ columns forming a basis of the column space. If two rows agree
on those columns, they agree on every linear combination of them and hence
on all columns. Restriction to the chosen columns therefore injects the
set of distinct rows into $S^r$, which has $s^r$ elements.
\end{proof}

Let $\mathcal A=\{0,1\}\times\{0,1\}$ and let
$S=\{0\}\cup\{z(x):x\in\{0,1\}^N\}$, so $|S|\leq m+1$.
For a prefix $p\in\mathcal A^k$, define the residual function
$r_p:\mathcal A^{N-k}\to S$ by $r_p(t)=f(pt)$, where juxtaposition is
concatenation. Let $Q_k$ be the set of distinct residual functions at cut
$k$. These are precisely the distinct rows of $M_k$. Lemma~\ref{lem:rows}
gives
\begin{align}
 |Q_k|\leq |S|^{\operatorname{rank}M_k}\leq W.
 \label{eq:width}
\end{align}
There is a well-defined transition
\begin{align}
 \delta_j:Q_{j-1}\times\mathcal A&\longrightarrow Q_j,&
 \delta_j(r_p,c)&=r_{pc}.
 \label{eq:transition}
\end{align}
Indeed, $r_p=r_{p'}$ implies $f(pct)=f(p'ct)$ for every suffix $t$.
The initial state $q_0$ is $r_{\varnothing}$.
Each terminal state $q\in Q_N$ is a scalar $\omega(q)\in S$.
The transitions and terminal values exactly recover $f$.

For matrix notation, embed each $Q_k$ in a common set of $W$ labels.
Define transitions arbitrarily on unused labels, always into the next
$Q_j$, and assign unused terminal labels weight zero. Actual trajectories
stay in $Q_j$, so this extension does not change the coefficients.
The zero residual, when present, is a rejecting state. No reversible
transition assumption is made: the classical computation will retain and
then erase its intermediate values.

\section{Recovering the output string}

Fix an input string $x=(x_1,\ldots,x_N)$. A terminal state is accepting
when its weight is nonzero. Let $t(q)$ be its accepting indicator.
For a state $q$ at cut $j$, let $s_j(q)$ indicate that some output suffix
completes the fixed input suffix $x_{j+1},\ldots,x_N$ to an accepting path.
Then
\begin{align}
 s_N(q)&=t(q),\\
 s_{j-1}(q)&=\bigvee_{a\in\{0,1\}}
       s_j\bigl(\delta_j(q,(a,x_j))\bigr).
 \label{eq:suffix}
\end{align}
For fixed $x_j$, define the Boolean matrix
\begin{align}
 E_j(q,q')&=\bigvee_{a\in\{0,1\}}
       [\delta_j(q,(a,x_j))=q'].
 \label{eq:relation}
\end{align}
Multiplication of Boolean matrices uses disjunction in place of addition
and conjunction in place of multiplication. Thus
$s_j=E_{j+1}\cdots E_Nt$.
All suffix matrices can be obtained simultaneously by products on a
balanced binary tree: first compute the product in each subtree; then
propagate the product strictly to its right from the root to the leaves.
This uses $O(N)$ matrix multiplications in $O(\log(N+1))$ stages.
Associativity suffices; commutativity is not needed.

For each state $q$ define $a_j(q)$ to be the least bit $a$ for which
$s_j(\delta_j(q,(a,x_j)))=1$. If no such bit exists, set $a_j(q)=0$.
Define the state transformation
\begin{align}
 h_j(q)&=\delta_j(q,(a_j(q),x_j)).
 \label{eq:greedy}
\end{align}
These tables depend only on $x_j$ and the suffix indicators, so every table
can be computed in parallel after (\ref{eq:suffix}).
All successive states are then obtained by a second computation of
associative products, now composing functions:
\begin{align}
 q_j&=(h_j\circ\cdots\circ h_1)(q_0),&
 y_j&=a_j(q_{j-1}).
 \label{eq:output}
\end{align}
Function composition may be represented by Boolean matrix multiplication,
with one nonzero entry in each row.

There is an accepting path for every input $x$, because $U$ is monomial.
Induction in (\ref{eq:greedy}) shows that the chosen prefix always admits
an accepting completion. Hence (\ref{eq:output}) ends at an accepting
terminal state. Such a path is unique for fixed $x$, again by monomiality.
It follows that $y=F(x)$ and $\omega(q_N)=z(x)$.
Uniqueness is needed only along the path from $q_0$; arbitrary states may
have several accepted completions, for which the least-bit rule is harmless.

For the inverse permutation, use $f^{\mathrm{op}}(a,b)=f(b,a)$ and ignore
its nonzero weights. Swapping the two indices at each site merely permutes
the rows and columns of every cut matrix, so its ranks remain at most $D$
and its coefficient alphabet is still $S$. Its unique accepted output for
input $y$ is $F^{-1}(y)$. Complex conjugation is unnecessary for recovering
the inverse permutation; phases are handled separately.

\section{From the Boolean computation to a quantum circuit}

A Boolean product of two $W$ by $W$ matrices requires $W^3$ conjunctions
and $O(W^3)$ binary disjunctions. A balanced disjunction tree has depth
$O(\log(W+1))$. Each matrix entry is used by at most $W$ conjunctions;
copying it to those uses takes $O(W)$ controlled-NOT gates and logarithmic
depth. Thus the computation has bounded fan-in and bounded fan-out after
copies are supplied, with size $O(W^3)$ and depth $O(\log(W+1))$ per product.
The tables $E_j,h_j,a_j$ and the evaluations in (\ref{eq:output}) use no
larger resources. An input bit $x_j$ is copied to its $O(W^2)$ table uses
with a binary tree of controlled-NOT gates.

The two balanced-tree computations therefore form a Boolean circuit for
$F$ of size $O(NW^3)$ and depth
$O(\log(N+1)\log(W+1))$, including copies. The same bounds hold for
$F^{-1}$. This is a circuit with fixed transition tables, not an algorithm
that reads a full $2^N$ by $2^N$ matrix during execution.

To make this circuit reversible, record each Boolean gate's output on a
fresh auxiliary bit. An AND is a Toffoli gate; an OR uses the identity
$a\lor b=a\oplus b\oplus ab$ and hence two controlled-NOT gates and one
Toffoli gate; a NOT and a copy have immediate reversible implementations.
Disjoint Boolean gates, with their previously supplied copies, can still
be applied in parallel. A Toffoli has an exact decomposition into a constant
number of one- and two-qubit gates \cite{Barenco95}.
Copy the desired output to a separate register and reverse the computation,
as in Bennett's construction \cite{Bennett73}. This gives
\begin{align}
 C_F:\ |x\rangle|b\rangle|0\rangle
 &\longmapsto |x\rangle|b\oplus F(x)\rangle|0\rangle
 \label{eq:reversible}
\end{align}
with a constant factor increase in size and depth and $O(NW^3)$ auxiliary
bits. The inverse-permutation circuit gives $C_{F^{-1}}$ in the same form.
The copies correlate computational-basis labels coherently; this reasoning
does not require, or assert, copying an arbitrary unknown quantum state.

For a monomial unitary, the computation of $F$ also supplies the terminal
state $q_N$. Encode it with $W$ bits, exactly one of which is one. For each
terminal label of nonzero weight, apply the one-qubit phase
$\operatorname{diag}(1,\omega(q))$ to its label bit; use the identity for
zero-weight labels. On the actual trajectory this multiplies the state by
$z(x)$. These gates act on separate bits and can be applied in one layer.
Insert this layer before reversing the computation that erases intermediate
values. Starting with a zero output register gives
\begin{align}
 |x\rangle|0\rangle|0\rangle
 &\longmapsto z(x)|x\rangle|F(x)\rangle|0\rangle\\
 &\longmapsto z(x)|F(x)\rangle|x\rangle|0\rangle\\
 &\longmapsto z(x)|F(x)\rangle|0\rangle|0\rangle.
 \label{eq:clean}
\end{align}
The second step swaps the two data registers, and the third applies
$C_{F^{-1}}$. The procedure acts correctly on every basis state and therefore
on every superposition. It proves Theorem~\ref{thm:monomial}.
Corollary~\ref{cor:permutation} takes $m=1$ and $W=2^D$.

\section{Reduction of arbitrary monomial MPUs to diagonal MPUs}

The permutation part can be separated even when the phases take arbitrarily
many values.

\begin{theorem}\label{thm:diagonal-reduction}
Let $U$ be any monomial MPU of bond dimension at most $D$.
Then $U=P\Phi$, where $P$ is a permutation MPU of bond dimension at most
$D^2$ and $\Phi$ is a diagonal MPU of bond dimension at most $D^3$.
The permutation $P$ has an exact circuit of $O(N8^{D^2})$ gates and auxiliary
qubits, and depth $O(D^2\log(N+1))$, with arbitrary two-qubit interactions.
\end{theorem}

\begin{proof}
Put $P_{yx}=|U_{yx}|^2$. Since $U|x\rangle=z(x)|F(x)\rangle$ and
$|z(x)|=1$, these entries are precisely $[y=F(x)]$.
If the local matrices of $U$ are $A_j(a,b)$, the coefficientwise squared
modulus has local matrices
\begin{align}
 B_j(a,b)&=A_j(a,b)\otimes\overline{A_j(a,b)},
 \label{eq:support-tensor}
\end{align}
and similarly tensored boundary vectors. The identity
$(A\otimes\overline A)(B\otimes\overline B)
=AB\otimes\overline{AB}$ proves that these tensors represent $P$ with
bond dimension at most $D^2$. Corollary~\ref{cor:permutation} implements it.
Now $\Phi=P^\dagger U$ acts as
$\Phi|x\rangle=z(x)|x\rangle$ and is therefore diagonal.
The product of two MPOs is obtained by summing the shared physical index
and tensoring their bond indices. Its bond dimension is at most the product
of the two bond dimensions, here $D^2D=D^3$.
\end{proof}

Thus polynomial circuit size for arbitrary diagonal MPUs of fixed bond
dimension would imply it for arbitrary monomial MPUs of fixed bond dimension.
The reverse implication is immediate because diagonal unitaries are
monomial. This reduction concerns circuit size with auxiliary qubits and
arbitrary two-qubit interactions; it imposes no finite phase alphabet.
It does not resolve the diagonal case with arbitrary phases.

\section{Construction of the finite tables and formalization}

The residual tables need not be found by enumerating all prefixes and
suffixes. Suppose an MPO representation is given by matrices $A_j(c)$,
left boundary $l$, and right boundary $r$. Starting with $V_N=\operatorname{span}\{r\}$,
compute backwards
\begin{align}
 V_{j-1}&=\operatorname{span}\{A_j(c)v:c\in\mathcal A,\ v\in V_j\}.
 \label{eq:suffix-span}
\end{align}
Keep a basis of each $V_j$, of dimension at most $D$.
A prefix has bond row $v_p=lA_1(c_1)\cdots A_j(c_j)$.
Two prefixes have the same residual if and only if their bond rows agree
on the selected basis of $V_j$. Starting from $l$, extend every retained
representative by all four local letters and retain one representative of
each distinct restricted row. Lemma~\ref{lem:rows} bounds the number retained
by $W$ at every step. Elementary linear algebra and exact equality tests
therefore construct the finite tables with $O(N\,\mathrm{poly}(D,W))$ field
operations. This is an exact arithmetic statement. No bound on bit complexity
for arbitrary complex tensor entries is claimed. The circuit existence and
gate bounds do not require a numerical separation between those entries.

The script \texttt{scripts/check\_mpu\_finite\_alphabet\_recovery.py}
checks output recovery and inverse erasure on every basis input of 424
examples: all 384 signed two-qubit permutations, 36 three- and four-qubit
finite-phase examples, and four multi-controlled NOT gates. Ranks and
comparisons are exact. Its enumeration of small residual tables verifies
the reconstruction, not compiled quantum gates or the construction cost.

The module \texttt{TNLean/MPS/MPU/FiniteAlphabetRows.lean} proves
Lemma~\ref{lem:rows}, the separating-column assertion, and the rank-bound,
finite-bond, and zero-one consequences. The remaining arguments in this note
are mathematical proofs rather than Lean theorems.

For arbitrary complex coefficients, the set of residual functions can grow
exponentially even when its span has dimension one, as for a product of
one-site phase gates. Such gates have efficient circuits, so this failure
of the finite-state bound gives no complexity lower bound. General
nonuniform MPUs require further structure beyond finite coefficient values.

\begin{thebibliography}{9}
\bibitem{STC25}
G.~Styliaris, R.~Trivedi, and J.~I.~Cirac,
\emph{Quantum Circuits for Matrix-Product Unitaries},
Physical Review Letters \textbf{135}, 260602 (2025),
\href{https://arxiv.org/abs/2508.08160}{arXiv:2508.08160v2}.
\bibitem{Bennett73}
C.~H.~Bennett, \emph{Logical Reversibility of Computation},
IBM Journal of Research and Development \textbf{17}, 525--532 (1973),
\href{https://doi.org/10.1147/rd.176.0525}{doi:10.1147/rd.176.0525}.
\bibitem{Barenco95}
A.~Barenco et al., \emph{Elementary gates for quantum computation},
Physical Review A \textbf{52}, 3457--3467 (1995),
\href{https://arxiv.org/abs/quant-ph/9503016}{arXiv:quant-ph/9503016}.
\end{thebibliography}
\end{document}
